% !TeX root = ../notes.tex
% Proposed prose version aligned to 06-choose-plan.tex.

\chapter{Choosing the next experiment}

\begin{notebox}
  \keypoint{What do you need to learn from this chapter?}
  Reaction optimisation is an experiment-selection problem defined by variables, objectives, constraints, available observations, and uncertainty.
  This chapter compares fixed and sequential strategies, then asks how chemical boundaries limit apparently rational recommendations.
\end{notebox}

\section{What is being optimised?}

An optimisation problem begins with a decision space.
Continuous variables can include temperature, time, concentration, pressure, flow rate, and reagent equivalents.
Categorical variables can include solvent, catalyst, ligand, base, or reactor architecture.
The feasible space is smaller than the formal set of combinations because solubility, equipment limits, safety, and chemical compatibility exclude some experiments.

The objective defines what counts as improvement.
A reaction with the highest analytical yield may not have the best isolated yield, selectivity, productivity, process mass intensity, cost, or hazard profile.
The objective function therefore encodes priorities rather than revealing a universal optimum.

\begin{conceptbox}
  \keypoint{An optimiser does not discover what ``best'' means.}
  It searches according to the objective and constraints supplied by the chemist.
\end{conceptbox}

\section{One variable at a time and hidden interactions}

One-variable-at-a-time optimisation changes one condition while holding the others fixed.
The procedure is intuitive, produces plots that are easy to interpret locally, and can be practical when experiments are expensive or difficult to reconfigure.
Its weakness is that the result depends on the starting point and on the order in which variables are examined.

An interaction occurs when the effect of one variable depends on the level of another.
For example, increasing temperature may improve conversion with one base but accelerate decomposition with another.
An OVAT experiment performed at only one base identity cannot reveal the full interaction.
The resulting path can stop at a locally favourable setting while missing a better region elsewhere in the response surface.

\section{Factorial thinking, design of experiments, and response surfaces}

Factorial designs vary several factors together so that main effects and interactions can be estimated.
A full factorial design measures every combination of the selected levels, while a fractional design uses a subset chosen to preserve specified information.
The important chemical question is not whether every combination is measured, but which effects become distinguishable within the experimental budget.

Two-level designs approximate the response locally as a plane.
Centre points and response-surface designs provide evidence for curvature when a linear local description is inadequate.
The fitted surface remains empirical, but it can identify influential variables, interactions, and promising regions more efficiently than a long sequence of uncoordinated tests.

\section{Sequential experimentation}

A fixed design is chosen before its outcomes are known.
A sequential campaign uses each completed round to update the description of the response and to choose subsequent experiments.
This adaptivity can concentrate measurements near promising, uncertain, or chemically informative regions.

Sequential does not necessarily mean one experiment at a time.
A batch can be selected when reactions run in parallel or when instrument scheduling dominates the cost.
The price of batching is that experiments within the batch are chosen before their neighbours have returned results.
A stopping rule is also required, such as reaching a target, exhausting the budget, or observing little expected improvement.

\section{Prediction with uncertainty}

A surrogate model approximates the expensive experimental response using the observations collected so far.
Gaussian-process models are useful conceptually because they return both a mean prediction and an uncertainty estimate over the search space.
The uncertainty is generally smaller near informative observations and larger in poorly sampled regions.

This uncertainty is conditional on the chosen model and representation.
It does not guarantee that a hidden mechanism, decomposition pathway, or analytical failure has been considered.
Experimental noise, model uncertainty, and ignorance of missing chemistry should not be collapsed into one unqualified confidence interval.

\section{Bayesian optimisation}

Bayesian optimisation combines a surrogate model with an acquisition function.
The acquisition function ranks candidate experiments by balancing predicted performance and uncertainty.
Exploitation selects a condition expected to perform well, while exploration selects a condition that may improve the model or reveal an unsampled region.
Expected improvement is one way to combine these motives into a single ranking.

Shields and co-workers benchmarked Bayesian reaction optimisation across several reaction datasets and compared machine recommendations with human decisions.\mcite{Shields2021}
The study shows why optimisation quality should be evaluated over complete trajectories rather than by displaying only the final successful condition.
A method that repeatedly reaches good regions with fewer experiments is more useful than one that occasionally finds an excellent point by chance.

\section{Multiple objectives and constraints}

Chemical optimisation often has several objectives.
Yield and selectivity may conflict with productivity, solvent use, energy demand, cost, or hazard.
A solution is dominated when another tested or predicted solution is at least as good in every objective and better in at least one.
The non-dominated solutions form a Pareto front.

The Pareto front does not choose among the trade-offs.
It exposes the consequences of prioritising one objective over another.
A chemist or process team must still decide whether a small gain in yield justifies additional waste, cost, or risk.
Unsafe or technically impossible regions should be represented as constraints rather than rewarded merely because the model is uncertain there.

\section{Pidko--Kalikadien: why model-guided choice can fail}

\subsection{Experimental-data and domain warning}

Kalikadien and co-workers generated 3552 measurements for rhodium-catalysed asymmetric hydrogenation across catalyst, substrate, solvent, pressure, temperature, and time choices.\mcite{Kalikadien2024}
The dataset contained many more negative or poor outcomes than the corresponding literature record.
This difference demonstrates how publication and selection practices alter the response distribution seen by a model.

The models described behaviour inside the sampled campaign more successfully than they generalised to chemistry outside it.
Expensive calculated descriptors did not automatically remove the limitation.
The case therefore changes the optimisation question from ``which algorithm should choose next?'' to ``what evidence shows that the candidate belongs to a domain where model-guided choice is meaningful?''

\subsection{Representation and conformational warning}

Baidun and co-workers examined descriptors for rhodium asymmetric-hydrogenation complexes using generic precatalyst and substrate-specific structures.\mcite{Baidun2024}
Substrate coordination altered the number, energy, and geometry of accessible conformers.
The associated steric and electronic descriptors therefore depended on which chemical complex and conformational ensemble was represented.

A virtual screen built from one convenient geometry can look numerically precise while omitting flexibility induced by the reacting substrate.
The lesson is not that every optimisation requires the most expensive calculation.
The lesson is that representation error must be considered before uncertainty from the fitted model is used to rank experiments.

\section{Optimisation across chemical boundaries}

A surrogate model usually assumes that the response changes smoothly enough to learn from nearby experiments.
A mechanistic boundary can violate that assumption when a catalyst changes ligation state, a different pathway becomes dominant, or a new phase appears.
The optimiser may cross the boundary and report improved conditions without recognising that the chemical regime has changed.

Abrupt residuals, changes in selectivity pattern, inconsistent analytical signals, or loss of model calibration should therefore trigger chemical investigation.
Orthogonal analysis and targeted mechanistic experiments can be more valuable at that point than another round of numerical optimisation.
Finding better conditions is not automatically the same as understanding why they are better.

\section{Covering space versus choosing the next experiment}

Chapter~4 selected representative molecules to cover a defined chemical space.
That operation is useful when an initial set should span a representation without reference to a measured response.
Chapter~6 instead selects an experiment because it serves an objective, reduces relevant uncertainty, or is expected to improve the current state of knowledge.
A representative point may be a sensible initial experiment, but coverage alone does not determine the best next experiment after outcomes are available.

\section{Summary}

Optimisation requires explicit variables, objectives, constraints, budgets, and stopping rules.
Factorial designs reveal interactions that OVAT can miss, while sequential designs allow observations to change later choices.
Bayesian optimisation combines prediction and uncertainty through an acquisition function.
Multiobjective optimisation identifies trade-offs rather than one universal optimum.
Dataset bias, representation choice, conformational dependence, and mechanistic boundaries limit the trust that should be placed in model-guided recommendations.
